\section{SRS загвар -
Nogoolin}\label{doc3-srs-ux437ux430ux433ux432ux430ux440---nogoolin}

\textbf{M2 / US-2.1 · Тэнгис · 2026-09-21 · v0.2 · SRS загвар}

\subsection{Загварыг ашиглах
нь}\label{doc3-ux437ux430ux433ux432ux430ux440ux44bux433-ux430ux448ux438ux433ux43bux430ux445-ux43dux44c}

IEEE 830-1998-ийн агуулгын бүтцийг ISO/IEC/IEEE 29148:2018-ийн шаардлага
тодорхойлох, шалгах, мөрдөх зарчимтай хослуулна. Доорх
\emph{{[}\ldots{]}} хэсэгт M3-аас агуулга нэмнэ. M2-д таван жишээ
шаардлагыг тусад нь бичсэн; энэ нь бүрэн SRS биш.

\textbf{Баримтын бүртгэл:} төсөл Nogoolin; зохиогч Тэнгис; хувилбар
v0.2; төлөв draft; хянагч \emph{{[}нэр{]}}; баталсан огноо
\emph{{[}огноо{]}}.

\subsubsection{1. Оршил}\label{doc3-ux43eux440ux448ux438ux43b}

\textbf{1.1 Зорилго.} \emph{{[}Энэ SRS ямар шийдвэр, хэрэгжилт,
шалгалтад ашиглагдах вэ?{]}} Гол уншигч нь хөгжүүлэгч Тэнгис; хоёрдогч
уншигч нь хичээлийн багш.

\textbf{1.2 Хамрах хүрээ.}
\href{../../requirements/scope-charter.md}{Scope Charter}-ийг мөрдөнө.
Эхлэх хүрээ: inquiry, wishlist, хувийн өгөгдлийн тусгаарлалт, hero
fallback, API тестийн хаалт. \emph{{[}Системийн хил ба амжилтын
хэмжүүр.{]}}

\textbf{1.3 Нэр томьёо.} \emph{{[}Нэр - нэг утгатай тайлбар.{]}} Эхлэл:
inquiry = бүтээгдэхүүний асуулга; wishlist = дараа үзэхээр хадгалсан
жагсаалт; FR = үйлдлийн шаардлага; NFR = чанарын шаардлага.

\textbf{1.4 Эх сурвалж.} \emph{{[}Баримтын нэр, хувилбар/огноо, хэсэг,
холбоос.{]}} Эхлэл: \href{https://github.com/Tengis01/Nogoolin}{Nogoolin
public repository}, W2 handout, Lecture 02,
\href{../../phase-0/README.md}{Phase 0 эхүүд}, IEEE 830, ISO/IEC/IEEE
29148.

\textbf{1.5 Бүтцийн тайлбар.} §2-т орчин ба хэрэглэгч, §3-т шаардлага;
хавсралтуудад mapping, чанарын checklist, traceability байна.
Технологийн шийдвэрийг \href{../../architecture/sdd-template.md}{SDD
загварт} бөглөнө.

\subsubsection{2. Ерөнхий
тодорхойлолт}\label{doc3-ux435ux440ux4e9ux43dux445ux438ux439-ux442ux43eux434ux43eux440ux445ux43eux439ux43bux43eux43bux442}

\textbf{2.1 Бүтээгдэхүүний орчин.} \emph{{[}Системийн зориулалт, хил,
гаднын системтэй харилцах байдал.{]}} Эхлэл: Nogoolin нь шашин, зан
үйлийн бүтээгдэхүүний каталог.

\textbf{2.2 Гол үйлдлүүд.} \emph{{[}Хэрэглэгчийн зорилгоор бүлэглэсэн
үйлдлийн жагсаалт; FR ID холбоос.{]}} Эхлэх жишээ: асуулга илгээх FR-01,
бүтээгдэхүүн хадгалах FR-02.

\textbf{2.3 Хэрэглэгчийн онцлог.} \emph{{[}Дүр, зорилго, мэдлэг,
хүндрэл, W1 persona-ийн эх.{]}} Guest, customer, admin-ийг ялгана;
\href{wiki/persona-and-pivot.md}{persona нэмэлт}-ээс эхэлнэ.

\textbf{2.4 Хязгаарлалт.} \emph{{[}Бизнес, орчин, нөөцийн зайлшгүй
хязгаар ба үндэслэл.{]}} Эхлэл: ганцаар хөгжүүлэх, монгол интерфэйс,
хүргэлт идэвхгүй.

\textbf{2.5 Таамаг ба хамаарал.} \emph{{[}Таамаг, эзэн, шалгах хугацаа,
буруу бол нөлөө.{]}} Local DB болон тестийн өгөгдөл бэлэн эсэхийг
шалгана; cloud, asset, stakeholder approval-ийг бэлэн гэж таамаглахгүй.

\textbf{2.6 Дараа хийх хүрээ.} \emph{{[}Боломж, хойшлуулах шалтгаан,
зорилтот үе.{]}} Scope Charter-ийн postponed жагсаалтыг холбоно.
Хичээлийн W3 ба бүтээгдэхүүний Phase 3-ыг адилтгахгүй.

\subsubsection{3. Тодорхой
шаардлагууд}\label{doc3-ux442ux43eux434ux43eux440ux445ux43eux439-ux448ux430ux430ux440ux434ux43bux430ux433ux443ux443ux434}

\textbf{3.1 Гадаад интерфэйс.} \emph{{[}Оролт/гаралт, формат, алдаа,
харилцах тал.{]}} 3.1.1 хэрэглэгчийн, 3.1.2 төхөөрөмжийн, 3.1.3
программын, 3.1.4 холбооны интерфэйс гэж салгана. Төхөөрөмжийн
интерфэйсийн n-a үндэслэлийг mapping-д тэмдэглэв.

\textbf{3.2 Функциональ шаардлага.} \emph{{[}Доорх шаардлагын картаар FR
бүрийг бич.{]}} M2-ийн \href{../../requirements/requirements.md}{FR-01,
FR-02} нь жишээ; үлдсэнийг W3-д нэмнэ.

\textbf{3.3 Гүйцэтгэл.} \emph{{[}Ачаалал, өгөгдлийн хэмжээ, орчин,
хугацааны хэмжүүр, pass/fail босго.{]}} ``Хурдан'' гэх мэт хэмжүүргүй
өгүүлбэр хэрэглэхгүй.

\textbf{3.4 Логик өгөгдөл.} \emph{{[}Хадгалах мэдээлэл, холбоо,
unique/ownership дүрэм, хадгалах хугацаа.{]}} Өгөгдлийн сангийн
бүтээгдэхүүн сонгох шийдвэрийг SDD-д бичнэ.

\textbf{3.5 Дизайнд тавих зайлшгүй хязгаар.} \emph{{[}Гаднаас тогтоосон
техникийн хязгаар, түүний эх ба үндэслэл.{]}} Өөрийн сонгосон
хэрэгжүүлэлтийн аргыг энд шаардлага болгож хуулж болохгүй.

\textbf{3.6 Чанарын шаардлага.} \emph{{[}Найдвартай ажиллагаа, бэлэн
байдал, аюулгүй байдал, засварлах боломж, зөөврийн байдал тус бүрийн
NFR.{]}} M2-ийн NFR-01\ldots03-ыг холбож, дутагдах хэсгийг W3-д нөхнө.

\textbf{3.7 Бусад шаардлага.} \emph{{[}Хувийн мэдээлэл, audit, хэрэглэх
бодлогын эх ба шалгуур.{]}} Мэдээлэл дутуу хэсгийг planned гэж
тэмдэглэнэ; тодорхойгүйг n-a гэж үзэхгүй.

\subsection{Хавсралт A. IEEE 830 section
mapping}\label{doc3-ux445ux430ux432ux441ux440ux430ux43bux442-a.-ieee-830-section-mapping}

\textbf{Төлөв:} \nolinkurl{complete} = M2-т шаардсан тухайн агуулга
бэлэн; бүтээгдэхүүн хэрэгжсэн гэсэн үг биш. \nolinkurl{planned (W…)} =
агуулгыг нөхөх үе. \nolinkurl{n-a (justified)} = хүрээнд хамаарахгүй,
нэг өгүүлбэр үндэслэлтэй. Мөр бүр яг нэг төлөвтэй.

\begin{longtable}[]{@{}
  >{\raggedright\arraybackslash}p{(\linewidth - 4\tabcolsep) * \real{0.2600}}
  >{\raggedright\arraybackslash}p{(\linewidth - 4\tabcolsep) * \real{0.5500}}
  >{\raggedright\arraybackslash}p{(\linewidth - 4\tabcolsep) * \real{0.1900}}@{}}
\toprule\noalign{}
\begin{minipage}[b]{\linewidth}\raggedright
Section
\end{minipage} & \begin{minipage}[b]{\linewidth}\raggedright
Content / агуулга, эх
\end{minipage} & \begin{minipage}[b]{\linewidth}\raggedright
Status
\end{minipage} \\
\midrule\noalign{}
\endhead
\bottomrule\noalign{}
\endlastfoot
1.1 Purpose & §1.1-ийн зорилго, уншигчийн хэрэгцээг эцэслэх & planned
(W03) \\
1.2 Scope & Нэг хуудас Scope Charter & complete \\
1.3 Definitions & Нэр томьёоны нэгдсэн тайлбар & planned (W03) \\
1.4 References & W2/лекц/стандарт/Phase 0 эхийн холбоос & complete \\
1.5 Overview & §1.5-ийн бүтцийн тайлбар & complete \\
2.1 Product perspective & Системийн хил, гаднын орчин; Vision & planned
(W03) \\
2.2 Product functions & FR-01/02-оос үйлдлийн хүрээг өргөжүүлэх &
planned (W03) \\
2.3 User characteristics & Persona нэмэлт, ярилцлагаар нягтлах & planned
(W03) \\
2.4 Constraints & Scope Charter, нөөц ба бизнесийн хязгаар & planned
(W03) \\
2.5 Assumptions/dependencies & Таамаг бүрийн эзэн, шалгах нөхцөл &
planned (W03) \\
2.6 Apportioning & Scope Charter-ийн postponed жагсаалт & complete \\
3.1.1 User interfaces & Хэрэглэгчийн урсгал, оролт, баталгаа/алдаа &
planned (W03) \\
3.1.2 Hardware interfaces & Энэ web/API хүрээнд тусгай төхөөрөмж,
сенсортой холболт байхгүй. & n-a (justified) \\
3.1.3 Software interfaces & Гаднын систем ба интерфэйсийн шаардлага;
нарийн API contract W5 & planned (W03) \\
3.1.4 Communications & Протокол, өгөгдлийн формат, хамгаалалтын
шаардлага & planned (W03) \\
3.2 Functional requirements & Хоёр FR жишээ; W3-д 15 FR болгох & planned
(W03) \\
3.3 Performance & Ачаалал, орчин, хэмжүүр, босго & planned (W03) \\
3.4 Logical database & Өгөгдөл ба бүрэн бүтэн байдлын дүрэм & planned
(W03) \\
3.5 Design constraints & Зайлшгүй хязгаар ба эх сурвалж & planned
(W03) \\
3.6.1 Reliability & Алдаа/сэргэлтийн шалгуур; NFR-02 эхлэл & planned
(W03) \\
3.6.2 Availability & Хэмжих хугацаа, бэлэн байдлын зорилт & planned
(W03) \\
3.6.3 Security & NFR-01; эрх ба хувийн өгөгдөл & planned (W03) \\
3.6.4 Maintainability & NFR-03; шаардлагыг W3, CI нотолгоог W9 & planned
(W03) \\
3.6.5 Portability & Дэмжих орчин ба шалгах нөхцөл & planned (W03) \\
3.7 Other requirements & Хувийн мэдээлэл, audit бодлого & planned
(W03) \\
Appendix B & Шаардлагын карт ба найман чанарын шалгуур & complete \\
Appendix C & 5 мөртэй, 8 баганат traceability; W1 холбоо & complete \\
\end{longtable}

\subsection{Хавсралт B. Шаардлагын карт ба чанарын
шалгалт}\label{doc3-ux445ux430ux432ux441ux440ux430ux43bux442-b.-ux448ux430ux430ux440ux434ux43bux430ux433ux44bux43d-ux43aux430ux440ux442-ux431ux430-ux447ux430ux43dux430ux440ux44bux43d-ux448ux430ux43bux433ux430ux43bux442}

Энэ картыг §3-ын шаардлага бүрд хуулж бөглөнө. ISO/IEC/IEEE 29148-ийн
шаардлагын шинж, атрибут, удирдлагын зарчмыг хэрэгжүүлэхэд ашиглана;
доорх бүх талбар стандартын үгчилсэн хуулбар биш.

\begin{longtable}[]{@{}
  >{\raggedright\arraybackslash}p{(\linewidth - 2\tabcolsep) * \real{0.2500}}
  >{\raggedright\arraybackslash}p{(\linewidth - 2\tabcolsep) * \real{0.7500}}@{}}
\toprule\noalign{}
\begin{minipage}[b]{\linewidth}\raggedright
Талбар
\end{minipage} & \begin{minipage}[b]{\linewidth}\raggedright
Загвар
\end{minipage} \\
\midrule\noalign{}
\endhead
\bottomrule\noalign{}
\endlastfoot
ID / төрөл / нэр & \emph{{[}FR-XX эсвэл NFR-XX{]} / {[}FR эсвэл NFR{]} /
{[}нэг үр дүнгийн нэр{]}} \\
Шаардлагын өгүүлбэр & \emph{{[}Нөхцөл{]} үед систем {[}ажиглагдах нэг үр
дүн{]}-г {[}must/shall/may{]} хангана.} \\
Priority & \emph{{[}must / shall / may{]}}; хичээлийн тэмдэглэгээг
мөрдөж \nolinkurl{should} хэрэглэхгүй \\
Source / үндэслэл & \emph{{[}Баримт/хэсэг/холбоос, persona pain point,
яагаад хэрэгтэй{]}} \\
Owner / source interview & \emph{{[}Нэр{]} / {[}ярилцсан хүн, огноо,
тэмдэглэлийн холбоос; хийгдээгүй бол хийгдээгүй{]}} \\
Acceptance criteria & \emph{Given {[}урьдчилсан нөхцөл{]}, When
{[}үйлдэл{]}, Then {[}ажиглах үр дүн ба pass/fail босго{]}.} \\
Verification & \emph{{[}Test / Inspection / Analysis{]} + {[}шалгалтын
ID, алхам, орчин, нотолгоо{]}} \\
Dependency / Risk & \emph{{[}Хамаарах ID/нөхцөл{]} / {[}эрсдэл{]}} \\
Status / Last-Reviewed & \emph{{[}Draft / reviewed / approved \ldots{]}
/ {[}YYYY-MM-DD{]}} \\
\end{longtable}

\subsubsection{W2-ын найман чанарын
шалгуур}\label{doc3-w2-ux44bux43d-ux43dux430ux439ux43cux430ux43d-ux447ux430ux43dux430ux440ux44bux43d-ux448ux430ux43bux433ux443ux443ux440}

\begin{longtable}[]{@{}
  >{\raggedright\arraybackslash}p{(\linewidth - 2\tabcolsep) * \real{0.2500}}
  >{\raggedright\arraybackslash}p{(\linewidth - 2\tabcolsep) * \real{0.7500}}@{}}
\toprule\noalign{}
\begin{minipage}[b]{\linewidth}\raggedright
Шалгуур
\end{minipage} & \begin{minipage}[b]{\linewidth}\raggedright
Хянах асуулт
\end{minipage} \\
\midrule\noalign{}
\endhead
\bottomrule\noalign{}
\endlastfoot
Correct & Хэрэгцээ нь эх сурвалжаар нотлогдож, stakeholder-аар
нягтлагдсан уу? \\
Unambiguous & Нэр томьёо, нөхцөл, үр дүнг нэг утгаар ойлгох уу? \\
Complete & Тухайн шаардлагын урьдчилсан нөхцөл, хэвийн/алдааны үр дүн
хангалттай юу? \\
Consistent & Бусад ID, Scope Charter-тай зөрчилдөх үү? \\
Ranked & Priority ил тод уу; өөрчлөгдөх магадлалыг шаардлагатай бол
тэмдэглэсэн үү? \\
Verifiable & Test/Inspection/Analysis-аар pass/fail-ийг тогтоож болох
уу? \\
Modifiable & Нэг ID нэг үүрэгтэй, давхар хуулбаргүй, засварлахад
ойлгомжтой юу? \\
Traceable & Эх хэрэгцээ → шаардлага → шалгалтын холбоо байгаа юу? \\
\end{longtable}

Энэ найман нэрийг Lecture 02-ын rubric-аар хэрэглэв. IEEE 830 §4.3-т мөн
эдгээр SRS чанарыг нэрлэдэг; ISO 29148-ийн бүх шинжийг зөвхөн энэ
наймаар хязгаарлаж ойлгохгүй.

\subsection{Хавсралт C. Traceability ба
өөрчлөлт}\label{doc3-ux445ux430ux432ux441ux440ux430ux43bux442-c.-traceability-ux431ux430-ux4e9ux4e9ux440ux447ux43bux4e9ux43bux442}

\href{../../requirements/traceability-matrix.md}{Traceability
Matrix}-ийн баганыг
\nolinkurl{ID | Source | Owner | Verification | Dependency | Risk | Status | Last-Reviewed}
гэсэн дарааллаар хадгална. Шаардлага бүр нэг мөртэй, дор хаяж гурван
багана бөглөсөн байна. M2-ийн таван жишээнд найман баганыг бөглөсөн.
Нэгээс доошгүй мөр W1 persona pain point-той холбоотой байна; төсөл
сольсон холбоог тайлбарлана.

Өөрчлөхдөө ID-г хадгалж, source, шалгуур, matrix, review огноог хамтад
нь шинэчилнэ. Тестийн нотолгоо байхгүй бол verified гэж тэмдэглэхгүй.
Бүрэн SRS (15 FR + 5 NFR) нь W3/M3-ын ажил.

\textbf{Эх:}
\href{Software_Project_Documentation_Week_02_Handout.pdf}{W2 handout},
\href{../lecture/Lecture_02\%20Why\%20Documentation\%20Documentation\%20Standards.pdf}{Lecture
02, слайд 4--8};
\href{https://www.math.uaa.alaska.edu/~afkjm/cs401/IEEE830.pdf}{IEEE
830-1998}, §4.3, §5;
\href{https://standards.ieee.org/ieee/29148/6937/}{ISO/IEC/IEEE
29148:2018}. Хичээлийн унших чиглэл: Bhatti бүлэг 2, х. 23--44 (хүрээ),
бүлэг 4, х. 67--81 (шаардлагын эх ба уншигч); Chinchilla бүлэг 5, х.
53--55 (хүлээлт, таамаг, одоогийн дуусах нөхцөл).
